\section{Expérience Professionnelle}


\cventry
{Responsable Sécurité des Systèmes d'Information \textbf{| RSSI}} % Titre
{Bourse des Valeurs Mobilières de l'Afrique Centrale \textbf{| BVMAC}} % Organisation
{CDI} % Contrat
{Depuis Oct. 2024} % Date
\begin{cvitems}
\item {\textbf{Gouvernance SI :} Élaboration et déploiement de la PSSI, cartographie complète des actifs et classification des données sensibles.}
\item {\textbf{Résilience :} Conception et mise en œuvre du PCA et du PRA ; pilotage des tests de reprise pour garantir la continuité opérationnelle.}
\item {\textbf{SOC/SIEM :} Déploiement et administration d'une solution SIEM (Wazuh), définition des cas d'usage, corrélation des alertes et réponse aux incidents 24/7.}
\item {\textbf{Pentest :} Planification et réalisation de tests d'intrusion (Black/Grey/White Box) sur le SI, les applicatifs et les réseaux (BVMAC + sites de réplication).}
\item {\textbf{Supervision réseau :} Mise en place d'outils de monitoring centralisés couvrant le réseau BVMAC, les SDBS et le site de repli.}
\item {\textbf{Conformité :} Pilotage des actions de conformité RGPD, recommandations COBAC et guide d'hygiène ANSSI ; gestion des risques (ISO 27005, EBIOS RM).}
\item {\textbf{BYOD \& IAM :} Implémentation d'une politique BYOD et d'un cadre de gestion des accès (moindre privilège, MFA, révocation des droits).}
\item {\textbf{Sauvegarde :} Définition et déploiement d'une stratégie de sauvegarde et de réplication des données critiques de l'infrastructure financière.}
\end{cvitems}


\cventry
{Security Engineer \textbf{| Ingénieur Sécurité} — Full Remote} % Titre
{adorsys GmbH \& Co. KG \textbf{| Fintech, Allemagne}} % Organisation
{CDI} % Contrat
{Depuis Déc. 2023} % Date
\begin{cvitems}
\item {\textbf{Open Source :} Développement d'une bibliothèque Java pour SD-JWT (identité numérique décentralisée) — contribution au projet adorsys/sd-jwt sur GitHub.}
\item {\textbf{IDS/IPS :} Automatisation du déploiement de Snort 3 avec Wazuh sur Kubernetes pour la détection et la prévention d'intrusions réseau (NFV).}
\item {\textbf{EPP/XDR/SIEM :} Déploiement d'une plateforme unifiée de protection des endpoints et de corrélation des événements de sécurité.}
\item {\textbf{SMSI :} Mise en place et maintien d'un Système de Management de la Sécurité de l'Information (SMSI) conforme à l'ISO 27001.}
\item {\textbf{DevSecOps :} Industrialisation des déploiements (Helm, Terraform, Tekton CI/CD, ArgoCD) sur AWS et GCP — intégration de la sécurité dès la conception.}
\end{cvitems}

\cventry
{Ingénieur Cybersécurité \textbf{| Projet de Fin d'Études}} % Titre
{Port Autonome de Kribi \textbf{| PAK}} % Organisation
{PFE} % Contrat
{Fév. 2023 — Nov. 2023} % Date
\begin{cvitems}
\item {\textbf{SIEM :} Analyse, conception et déploiement de Wazuh pour la supervision et la corrélation des événements de sécurité sur l'infrastructure portuaire.}
\item {\textbf{Détection :} Développement de règles personnalisées et de décodeurs adaptés aux flux spécifiques du SI (logs réseau, systèmes, applicatifs).}
\item {\textbf{IDS/IPS :} Intégration de Snort pour la détection des menaces sur le périmètre réseau.}
\end{cvitems}

\cventry
{Ingénieur Sécurité Réseaux \textbf{| Stage}} % Titre
{Délégation Générale à la Sûreté Nationale \textbf{| DGSN}} % Organisation
{Stage} % Contrat
{2022} % Date
\begin{cvitems}
\item {\textbf{NAC :} Conception et déploiement d'un contrôle d'accès réseau (Network Access Control) pour sécuriser les accès au SI.}
\item {\textbf{Audit :} Cartographie des actifs SI, analyse des risques et mise en place d'une politique de surveillance du réseau.}
\end{cvitems}
